% Reconstructed after the workspace reset; proofs and freshly regenerated
% exact certificates are in r275/, r216/, and r280/.
\begin{theorem}[Dimension-80 ring rigidity]\label{thm:m3-dimension80}
Every complex fusion category of type $(1,5;5,3)$ has the based ring of
$\operatorname{Rep}((C_2)^4\rtimes C_5)$ for the fixed-point-free action.
Its maximum fusion multiplicity is two. Hence $R8\_3\_275$,
$R8\_3\_276$, $R8\_3\_344$, and $R8\_3\_345$ are not categorifiable.
\end{theorem}
\begin{proof}
The category is solvable by \cite[Theorem 1.6]{ENOSolvable}. Use its
cyclic-prime extension/equivariantization sequence
\cite[Definition 1.2(ii), Proposition 4.4]{ENOSolvable}. A final
extension would have neutral dimension $40$ or $16$, neither of the
form $a+25b$ with $a\in\{1,5\}$ and $0\le b\le3$.
Thus $\mathcal C=\mathcal D^{C_5}$, since its invertibles force the
equivariantization prime to be five, and $\operatorname{FPdim}\mathcal D=16$.

A fixed base simple of dimension $d$ contributes five simples of
dimension $d$, and a free orbit contributes one of dimension $5d$;
cyclic stabilizers have trivial projective multipliers. Base dimensions
are rational algebraic integers, hence integral. A fixed dimension-five
simple is impossible since $25>16$, and all free-orbit simples must be
invertible. Thus $\mathcal D=\mathrm{Vec}_H^\omega$, $|H|=16$, with
only the unit fixed by $C_5$.
Every invariant subset containing the identity has size one modulo five.
The nontrivial characteristic center of the $2$-group $H$ therefore has
order sixteen; then its characteristic subgroup $H[2]$ also has order
sixteen. Hence $H=(C_2)^4$. The order-five action has irreducible minimal
polynomial $\Phi_5(t)=t^4+t^3+t^2+t+1$ over $\mathbb F_2$, so it is
unique up to conjugacy.

The five unit characters act trivially on the three free-orbit simples
$X_O$, and $X_O^*=X_O$. Reciprocity gives
$[X_OX_P:\chi]=\delta_{O,P}$ for every unit character $\chi$.
For a nonzero orbit $Q$, the coefficient at $X_Q$ is the number of
pairs $(x,y)\in O\times P$ with $x+y=z$, for any fixed $z\in Q$.
This follows by forgetting equivariance and is independent of the
associator and action coherence. The exact three-orbit convolution
certificate has maximum two, whereas the four source tensors have
maximum three.
\end{proof}

\begin{theorem}[A dimension-216 solvability obstruction]\label{thm:m3-dimension216}
No complex fusion category has type $(1,2;2,1;3,2;8,3)$.
In particular, $R8\_3\_810$, $R8\_3\_812$, $R8\_3\_2492$,
$R8\_3\_2494$, $R8\_3\_2620$, and $R8\_3\_2622$ are not categorifiable.
\end{theorem}
\begin{proof}
Use the same solvability sequence. A final extension would have neutral
dimension $108$ or $72$, neither expressible as $a+4b+9c+64d$ with
$a\in\{1,2\}$, $b\le1$, $c\le2$, $d\le3$. Thus
$\mathcal C=\mathcal D^{C_2}$, with base dimension $108$.
The cyclic orbit rule forces precisely the base types
\[
 A=(1,3;3,1;4,6),\qquad B=(1,3;3,1;4,2;8,1).
\]
Integrality descends by the orbit dimension formula and rational
algebraic integrality. Neither base type permits subcategory dimension
$54$ or $36$: enumerate $a+9b+16c+64d$, $a\in\{1,3\}$, within the
displayed bounds. Hence $\mathcal D=\mathcal E^{C_3}$, with base
dimension $36$. Type $B$ is impossible because its two dimension-four
simples would have to arise from fixed objects in multiples of three.
Type $A$ forces $\mathcal E$ of type $(1,4;4,2)$, and the actual
$C_3$ action cycles its three nonidentity invertibles. Their group is
therefore $V_4$.

For a dimension-four simple $X$, reciprocity gives
$16=|\operatorname{Stab}(X)|+4m$. The stabilizer order divides four,
so it is four. Thus all invertibles and a noninvertible belong to
$\mathcal E_{\rm ad}$, whose dimension is $20$ or $36$.
The ratio $\operatorname{FPdim}(\mathcal E)/
\operatorname{FPdim}(\mathcal E_{\rm ad})=|U(\mathcal E)|$ is integral;
therefore $U(\mathcal E)=1$.
The final solvability step is consequently a $C_2$ equivariantization,
giving a central Tannakian lift of one nontrivial invertible. Transport
it under the actual $C_3$ action to obtain bosonic central lifts of all
three nontrivial invertibles.

The kernel of $\operatorname{Inv}Z(\mathcal E)\to
\operatorname{Inv}\mathcal E$ is
$\operatorname{Aut}_{\otimes}(\operatorname{Id}_{\mathcal E})
=\operatorname{Hom}(U(\mathcal E),\mathbb C^\times)=1$.
Hence these lifts form a closed $V_4$, and its quadratic self-braiding
form is identically one. This pointed braided subcategory is Tannakian.
De-equivariantization gives $\mathcal E=\mathcal F^{V_4}$, where
$\operatorname{FPdim}\mathcal F=9$ and integrality again descends.
The integral prime-square theorem \cite{ENO,EGNO} makes
$\mathcal F=\mathrm{Vec}_H^\omega$, with $H=C_9$ or $C_3^2$.

An orbit of length two has cyclic stabilizer and produces forbidden
dimension-two simples. A nonidentity fixed point has stabilizer $V_4$;
its projective degrees are either $1,1,1,1$ or $2$. These would produce
four extra invertibles, beyond the four from the unit, or a forbidden
dimension-two simple. Thus all nonidentity elements of $H$ have free
$V_4$ orbits. Each nontrivial involution therefore acts fixed-point-freely.
An odd abelian group has only one fixed-point-free involution, namely
inversion: $(f-1)(f+1)=0$ and $\ker(f-1)=0$ imply $f=-1$.
Three nonidentity elements of $V_4$ cannot all act by inversion.
This contradiction proves the theorem.
\end{proof}

\begin{theorem}[A split-module obstruction in dimension 280]\label{thm:m3-dimension280}
The based rings $R8\_3\_1051$ and $R8\_3\_2415$ are not categorifiable
over $\mathbb C$.
\end{theorem}
\begin{proof}
The exact source certificate puts both rings in dimensions
$[1,1,5,5,6,8,8,8]$, with the nontrivial invertible $g$ swapping the
unit pair and the dimension-five pair and fixing the remaining four
objects $U_i$. For a self-dual dimension-five object $X$, its coefficient
block among the $U_i$ is
\[
 R=\begin{pmatrix}1&1&1&1\\1&0&2&1\\1&2&0&1\\1&1&1&1\end{pmatrix},
\]
and the source has the exact based-ring identity
$p(X)=0$, where $p(t)=t(t-5)(t-1)(t+1)(t+2)$.

The fixed dimension-six simple supplies a rank-one module category for
$\mathrm{Vec}_{C_2}^{\omega}$, so $\omega$ is a coboundary and the
regular algebra $A=1+g$ exists. In the category of right $A$-modules,
the two free $g$-orbits give two simples, and each fixed $U_i$ splits
into $U_{i,+},U_{i,-}$; the stabilizer multiplier vanishes since
$H^2(C_2,\mathbb C^\times)=0$.

The action of self-dual $X$ is symmetric. Induction is a module functor,
so its block between two split pairs has equal row and column sums
$R_{ij}$. That block is
$\left(\begin{smallmatrix}u&R_{ij}-u\\R_{ij}-u&u\end{smallmatrix}\right)$,
with $0\le u\le R_{ij}$. Coefficients connecting a free simple to the
two members of a split pair agree, since the free simple is induced
and induction sends a fixed object to the sum of its two summands.
Thus the four pair differences form an invariant subspace.

After independently interchanging the two summands of each pair, the
matrix on that subspace must be
\[
 S=\begin{pmatrix}a&1&1&1\\1&0&b&c\\1&b&0&d\\1&c&d&e\end{pmatrix},
 \qquad a,c,d,e\in\{-1,1\},\quad b\in\{-2,0,2\}.
\]
Every module action must satisfy the same based-ring identity, hence
$p(S)=0$. The certificate evaluates the full polynomial over the
integers for all $48$ possibilities and records a nonzero entry in
each case. It also verifies $p(N_X)=0$ in each source regular
representation, which is faithful. This gives the required contradiction
without a unitarity assumption.
\end{proof}
